\documentclass{resume}
\usepackage{zh_CN-Adobefonts_external}
\usepackage{linespacing_fix}
\usepackage{cite}
\usepackage{hyperref}
\hypersetup{
    colorlinks=true,
    linkcolor=cyan,
    filecolor=magenta,
    urlcolor=blue,
}

\begin{document}
\pagenumbering{gobble}

%***********个人信息**************
\MyName{崔兆诺}
\sepspace
\SimpleEntry{\textbf{籍贯：}山东省泰安市}
\SimpleEntry{\textbf{邮箱：}cuizhaonuo24@mails.ucas.ac.cn}
\SimpleEntry{\textbf{电话：}18325480212}
\SimpleEntry{\textbf{政治面貌：}群众}
\SimpleEntry{\textbf{通讯地址：}北京市朝阳区}

%************照片**************
%\yourphoto{0.13}  % 不需要照片

%***********教育背景**************
\section{教育背景}
\datedsubsection{\textbf{中国科学院大学}，生物与医药，\textit{硕士}}{2024.9 - 至今}

\datedsubsection{\textbf{山东农业大学}，计算机科学与技术，\textit{本科}}{2020.9 - 2024.6}


%***********项目经历**************
\section{项目经历}

\datedsubsection{\textbf{小麦小穗数高通量自动获取}}{2026.3 - 2026.7}
{工作内容：基于MaskDINO构建端到端小穗计数模型，实现麦穗分割与小穗定位一体化}
\Content
{ 设计嵌套式两级匈牙利匹配机制，外层匹配麦穗实例、内层匹配小穗点，实现分割与计数一体化建模；}
{ 提出框相对归一化坐标回归策略，将小穗点约束在所属麦穗框内，缓解跨穗漂移与背景误定位；}
{ 引入穗内k近邻距离约束，利用小穗沿穗轴规律排列的空间先验，增强密集小穗点匹配收敛性。}

\datedsubsection{\textbf{植物叶片病害半监督识别与鲁棒学习}}{2026.5 - 2026.7}
{工作内容：基于Mean Teacher半监督框架优化植物病害识别，解决标注稀缺与噪声标签问题}
\Content
{ 设计不确定性感知Mean Teacher半监督学习框架，通过学生-教师一致性训练利用少量标注数据提升模型泛化能力；}
{ 引入DWT频域特征增强，将低频结构与高频纹理信息融合到CNN特征中，增强早期病斑与细粒度病害表达；}
{ 设计监督对比学习与KL散度伪标签修正策略，优化特征分布并缓解噪声标签与误差累积问题。}

\datedsubsection{\textbf{密集麦穗实例分割与检测算法优化}}{2025.9 - 2026.2}
{工作内容：基于MaskDINO构建麦穗实例分割与检测模型，优化密集遮挡场景精度}
\Content
{ 设计多分支one-to-one/one-to-many匹配训练机制，通过解耦query学习增强密集目标监督信号；}
{ 提出Query-Level去重与位置感知匹配策略，从query去重和位置约束角度提升训练稳定性与预测质量；}
{ 在Swin-L MaskDINO强baseline上实例分割AP+3.4\%、检测AP+4.78\%，密集场景检测效果显著提升。}

\datedsubsection{\textbf{山东省软件设计比赛}，省一等奖}{2023.6 - 2023.9}
{工作内容：设计深度学习病虫害识别模型，开发智能问答与论坛等后端功能}
\Content
{ 设计并优化基于深度学习的病虫害图像识别模型，实现作物病害自动分类与防治知识精准推荐；}
{ 基于AC自动机实现病虫害知识库规则匹配智能问答系统，支撑常见病虫害问题自动应答；}
{ 搭建SpringBoot后端框架与RESTful API，实现论坛交流、知识查询等核心功能模块。}


%***********个人评价**************
\section{个人评价}

{CET-4/6通过，具有良好的英文文献阅读与学术写作能力，能够独立撰写技术报告与学术论文}

{精通Python、PyTorch与OpenCV，具备深度学习模型训练与优化的项目实践经验}

{获校三等奖学金，微视频创作比赛校三等奖、演讲比赛学院二等奖，综合素质全面发展}

{热爱篮球运动，曾作为学院篮球队队员参加院级比赛，具备良好的团队协作与沟通能力}

\end{document}
